\section{Fourier transforms and normalization}\label{app:A.3}

A function of position is transformed with a plus sign in the exponent,
\begin{equation}
  O(\vect{q}) = \int\dd^3r\;\ee^{+\ii\vect{q}\cdot\vect{r}}\,O(\vect{r}), \qquad
  O(\vect{r}) = \int\frac{\dd^3q}{(2\pi)^3}\;\ee^{-\ii\vect{q}\cdot\vect{r}}\,O(\vect{q}),
  \label{eq:A.3-space}
\end{equation}
and a function of time with the same sign,
\begin{equation}
  f(\omega) = \int\dd t\;\ee^{+\ii\omega t}f(t), \qquad
  f(t) = \int\frac{\dd\omega}{2\pi}\;\ee^{-\ii\omega t}f(\omega).
  \label{eq:A.3-time}
\end{equation}
The factor $(2\pi)^{-1}$ per dimension is in the inverse transform. Applied to the density of point electrons, $\nd_e(\vect{r}) = \sum_i\delta^3(\vect{r} - \vect{r}_i)$, \cref{eq:A.3-space} gives the density operator
\begin{equation}
  n_e(\vect{q}) = \sum_i \ee^{+\ii\vect{q}\cdot\vect{r}_i},
  \label{eq:A.3-ne}
\end{equation}
and applied to the dark-matter--electron potential it gives $\tilde V(\vect{q}) = \int\dd^3x\,\ee^{+\ii\vect{q}\cdot\vect{x}}V(\vect{x})$, the form used in \cref{eq:1.4-Vq}.

\subsection{Sign of the spatial transform}

The sign is fixed by the direction of $\vect{q}$, which is the momentum the dark matter gives to the target (\cref{app:A.4}). Take dark-matter plane waves $\langle\vect{r}_\chi|\vect{p}\rangle = \Vol^{-1/2}\ee^{\ii\vect{p}\cdot\vect{r}_\chi}$ in a box of volume $\Vol$, and a potential $V(\vect{r}_\chi - \vect{r}_e)$ between the dark matter and one electron. Shifting the integration variable to $\vect{x} = \vect{r}_\chi - \vect{r}_e$,
\begin{equation}
  \langle\vect{p} - \vect{q}|V(\vect{r}_\chi - \vect{r}_e)|\vect{p}\rangle
  = \frac{1}{\Vol}\int\dd^3r_\chi\,\ee^{\ii\vect{q}\cdot\vect{r}_\chi}V(\vect{r}_\chi - \vect{r}_e)
  = \frac{1}{\Vol}\,\tilde V(\vect{q})\,\ee^{+\ii\vect{q}\cdot\vect{r}_e} .
  \label{eq:A.3-vertex}
\end{equation}
With \cref{eq:A.3-space}, the transform of the potential and the operator left on the target both carry $+\ii\vect{q}\cdot\vect{r}$, and summed over electrons the target operator is $n_e(\vect{q})$ of \cref{eq:A.3-ne}. With the opposite spatial convention, $\tilde V$ would be evaluated at $-\vect{q}$ and the target operator would be $n_e(-\vect{q})$. Single-particle form factors follow the same rule: $f_{i\to f}(\vect{q}) = \langle f|\ee^{\ii\vect{q}\cdot\vect{r}}|i\rangle$~\cite[Eq.~(3.10)]{essig2016direct}. Ref.~\cite[Eq.~(S.4)]{hochberg2021determining} uses the other spatial convention with the same direction of $\vect{q}$, and writes the density operator as $\sum_i\ee^{-\ii\vect{q}\cdot\vect{r}_i}$; Ref.~\cite[Eq.~(4)]{trickle2022effective} keeps $\ee^{+\ii\vect{q}\cdot\vect{x}}$ on the target by writing the potential's transform at $-\vect{q}$. \Cref{app:A.4} gives the condition under which these differences matter.

The time convention fixes where retarded functions are analytic. A response that starts at $t = 0$ and oscillates at a transition energy $\Omega$ transforms as $-\ii\int_0^\infty\dd t\,\ee^{\ii(\omega - \Omega)t - 0^+t} = 1/(\omega - \Omega + \ii 0)$. Its pole lies in the lower half of the complex $\omega$ plane, so retarded functions are analytic for $\Imag\omega > 0$, and their denominators carry $\omega - (E_f - E_i) + \ii 0$ (\cref{sec:4.1}). A time-dependent perturbation of frequency $\omega > 0$ is $\propto\ee^{-\ii\omega t}$, the convention of quantum mechanics; engineering texts use $\ee^{+\ii\omega t}$ (written $\ee^{j\omega t}$), and their dielectric functions are the complex conjugates of ours, with $\Imag\varepsilon < 0$ for absorption. Here $\varepsilon = \varepsilon_1 + \ii\varepsilon_2$ with $\varepsilon_2 > 0$ for an absorbing medium at $\omega > 0$.

\subsection{Box normalization}

States are normalized in a box of volume $\Vol$ with periodic boundary conditions, and the limit $\Vol \to\infty$ is taken at the end. Allowed wavevectors form a grid of density $\Vol/(2\pi)^3$, so sums and Kronecker deltas become integrals and delta functions by
\begin{equation}
  \sum_{\vect{q}} \;\longleftrightarrow\; \Vol\!\int\frac{\dd^3q}{(2\pi)^3}, \qquad
  \delta_{\vect{q}\vect{q}'} \;\longleftrightarrow\; \frac{(2\pi)^3}{\Vol}\,\delta^3(\vect{q} - \vect{q}').
  \label{eq:A.3-box}
\end{equation}
The second follows from the first, since $\sum_{\vect{q}}\delta_{\vect{q}\vect{q}'} = 1$. The box contains $N$ unit cells of volume $\Omega$, so $\Vol = N\Omega$, and a uniform grid of crystal momenta in the first Brillouin zone has $N$ points (\cref{sec:3.1}). Sums over the zone become
\begin{equation}
  \frac{1}{N}\sum_{\vect{k}\in\mathrm{BZ}} \;\longleftrightarrow\; \Omega\int_{\mathrm{BZ}}\frac{\dd^3k}{(2\pi)^3}, \qquad
  \langle\!\langle \cdots \rangle\!\rangle \equiv \frac{1}{N^2}\sum_{\vect{k},\vect{Q}} (\cdots) = \Omega^2\int_{\mathrm{BZ}}\frac{\dd^3k\,\dd^3Q}{(2\pi)^6}\,(\cdots),
  \label{eq:A.3-BZ}
\end{equation}
where the double bracket is the zone average over an electron momentum $\vect{k}$ and a phonon momentum $\vect{Q}$, used for phonon-assisted processes (\cref{sec:6.5}).

Every intensive quantity must come out independent of $N$, and the powers of $N$ and $\Vol$ make that check mechanical. The structure factor of \cref{sec:4.2} has a prefactor $2\pi/\Vol$, an independent-particle sum $\sum_{\vect{k}}$ contributes a factor $N$, and the density vertex $\rho_{mn}(\vect{k},\vect{q})$ is of order one (\cref{app:A.5}). Their product is $2\pi N/\Vol = 2\pi/\Omega$, independent of the size of the crystal. A factor of $N$ or $\sqrt N$ that survives this count signals an error in a normalization. In the same way the rate per unit target mass contains $1/\rho_T$ and no $N$ (\cref{app:A.8}).

\subsection{Lattice-periodic functions}

A function with the periodicity of the lattice, such as the crystal potential or the periodic part of a Bloch state, has components only at reciprocal-lattice vectors. With \cref{eq:A.3-space} restricted to one cell,
\begin{equation}
  f_{\vect{G}} = \frac{1}{\Omega}\int_\Omega\dd^3r\;\ee^{+\ii\vect{G}\cdot\vect{r}}f(\vect{r}), \qquad
  f(\vect{r}) = \sum_{\vect{G}} f_{\vect{G}}\,\ee^{-\ii\vect{G}\cdot\vect{r}},
  \label{eq:A.3-periodic}
\end{equation}
and the transform over the whole crystal is
\begin{equation*}
  f(\vect{q}) = \Vol\sum_{\vect{G}}f_{\vect{G}}\,\delta_{\vect{q},\vect{G}} \;\longrightarrow\; (2\pi)^3\sum_{\vect{G}}f_{\vect{G}}\,\delta^3(\vect{q} - \vect{G}) ,
\end{equation*}
by \cref{eq:A.3-box}. Plane-wave codes expand in the opposite direction, with coefficients $c_{n\vect{k}}(\vect{G})$ of $\ee^{+\ii\vect{G}\cdot\vect{r}}$ in $u_{n\vect{k}}(\vect{r})$~\cite[Eq.~(3.14)]{essig2016direct}. A code coefficient $c(\vect{G})$ is therefore our component at $-\vect{G}$, and a formula that combines code coefficients with the components of \cref{eq:A.3-periodic} needs $\vect{G} \to -\vect{G}$ in one of them. For a real function, $f_{-\vect{G}} = f_{\vect{G}}^*$, and the distinction only conjugates; for a Bloch function it does not. Ref.~\cite[Appendix~A]{giustino2017electron} states the transforms between the direct and the reciprocal lattice that the electron--phonon literature uses.

\subsection{Crystal momenta and reciprocal-lattice matrices}

In a crystal, a momentum transfer $\vect{q}$ is split into a first-zone part and a reciprocal-lattice vector, $\vect{q} = \vect{q}_r + \vect{G}_q$ with $\vect{q}_r$ in the first Brillouin zone (\cref{sec:3.1}). Response functions are matrices in reciprocal-lattice vectors, such as $(\varepsilon^{-1})_{\vect{G}\vect{G}'}(\vect{q}_r,\omega)$: the transform of a function of two positions $\vect{r}$ and $\vect{r}'$ with $\ee^{+\ii(\vect{q}_r+\vect{G})\cdot\vect{r}}$ and $\ee^{-\ii(\vect{q}_r+\vect{G}')\cdot\vect{r}'}$. The book indexes such a matrix by the unreduced $\vect{q}$ at $\vect{G} = 0$:
\begin{equation}
  (\varepsilon^{-1})_{00}(\vect{q},\omega) \equiv (\varepsilon^{-1})_{\vect{G}_q\vect{G}_q}(\vect{q}_r,\omega),
  \label{eq:A.3-unreduced}
\end{equation}
and likewise for every reciprocal-lattice matrix. Both forms appear in the literature, often for the same object; Ref.~\cite[Eq.~(24)]{knapen2021dark} writes the second with an explicit selection $\delta_{\vect{k},|\vect{q}+\vect{G}|}$. The momentum transfer $\vect{q}$ itself is never folded: it is the momentum the dark matter loses, and the dark-matter kinematics of \cref{sec:1.3} depend on $|\vect{q}|$, not on $|\vect{q}_r|$.
